%!TEX root = ../../main.tex
\section{킬이 없는 전투}
\label{sec:eng-killless}

본 논문의 분석 대상은 킬을 기준으로 찾은 교전이다. 챔피언이 죽지 않는 전투는 이 정의로 잡히지 않는다. 상대의 기술을 소모시키고 물러나게 하는 것만으로도 전투에서 이길 수 있으므로, 그런 국면이 얼마나 흔한지와 잡아낸다 해도 결과를 잴 수 있는지를 따져 둔다.

킬 기반 검출기는 자신이 빠뜨린 것을 셀 수 없다. 그래서 킬 대신 근접을 세었다. 정의 상수를 추정한 코퍼스에서 시드를 고정해 뽑은 2만 경기의 5초 격자에서, 살아 있는 어떤 챔피언이 자신을 포함해 각 팀의 살아 있는 챔피언을 $k$명 이상 \constR{} 단위 안에 두는 구간을 근접 국면으로 보았다. 그 구간 안이나 양 끝 \constB{}초 안에 지도 어디서든 킬이 없으면 킬 없는 국면으로 세었다. 한타의 문턱인 $k = 4$에서 킬 없는 비율은 위치를 그리는 두 방식에 걸쳐 \killlessFourLo{}--\killlessFourHi{}\%이고, $k = 2$에서 \killlessTwoLo{}--\killlessTwoHi{}\%이다. Schubert 등 \cite{schubert2016encounter}이 Dota~2에서 보고한 비율은 \killlessSchubert{}\%였다. 분모는 근접 국면이고 근접이 곧 싸울 뜻은 아니므로, 이 비율은 빠뜨린 전투의 수가 아니라 근접 국면 가운데 킬이 없는 비율이다.

잡아낸다 해도 결과를 잴 수 없다. 본 논문의 결과는 예측 시점과 끝 시각을 첫 킬과 마지막 킬에서 정하므로, 킬이 없으면 두 시각이 모두 없다. 이전 연구의 교환가치 라벨도 같은 제약을 갖는다. 그 라벨은 시작 시각을 첫 킬에서 정하고 귀속 범위를 첫 킬 위치에서 정한다. 시각과 중심을 억지로 부여해도, 값이 매겨지는 사건이 없으면 점수는 0이 된다. 검출한 킬 없는 국면에 그 라벨을 적용하면 $k = 4$에서 23.3\%만이 값이 매겨지는 사건을 갖는다. 따라서 이 자료에서 킬 없는 전투는 찾는 단계와 재는 단계 모두에서 분석 대상 밖에 있다.
